\section{Estruturas de Dados}

\subsection{Fenwick Tree}

Suporta soma de prefixo e atualização pontual.

\textbf{Atualização:} $O(\log N)$

\textbf{Consulta de prefixo:} $O(\log N)$

\textbf{Soma em $[l,r]$:} $O(\log N)$

Usa índices de $1$ a $N$.

\begin{lstlisting}
int n;
vector<long long> bit;

void init(int size) {
    n = size;
    bit.assign(n + 1, 0);
}

void add(int i, long long x) {
    for (; i <= n; i += i & -i)
        bit[i] += x;
}

long long sum(int i) {
    long long res = 0;

    for (; i > 0; i -= i & -i)
        res += bit[i];

    return res;
}

long long query(int l, int r) {
    return sum(r) - sum(l - 1);
}
\end{lstlisting}

Uso:
\begin{lstlisting}
init(n);
add(i, x);
long long ans = query(l, r);
\end{lstlisting}


\subsection{Segment Tree}

Suporta consultas em intervalos e atualizações pontuais.

\textbf{Construção:} $O(N)$.

\textbf{Consulta:} $O(\log N)$.

\textbf{Atualização pontual:} $O(\log N)$.

\begin{lstlisting}
int n;
vector<long long> tree;

long long build(const vector<long long>& a, int l, int r, int p) {
    if (l == r)
        return tree[p] = a[l];

    int m = (l + r) >> 1;

    return tree[p] = build(a, l, m, p << 1) + build(a, m + 1, r, (p << 1) | 1);
}

long long query(int l, int r, int p, int ql, int qr) {
    if (r < ql || l > qr)
        return 0;

    if (l >= ql && r <= qr)
        return tree[p];

    int m = (l + r) >> 1;

    return query(l, m, p << 1, ql, qr) + query(m + 1, r, (p << 1) | 1, ql, qr);
}

long long update(int l, int r, int p, int idx, long long value) {
    if (r < idx || l > idx)
        return tree[p];

    if (l == r)
        return tree[p] = value;

    int m = (l + r) >> 1;

    return tree[p] = update(l, m, p << 1, idx, value) + update(m + 1, r, (p << 1) | 1, idx, value);
}
\end{lstlisting}

Inicialização:
\begin{lstlisting}
n = a.size();
tree.assign(4 * n, 0);
build(a, 0, n - 1, 1);
\end{lstlisting}

Uso:
\begin{lstlisting}
update(0, n - 1, 1, idx, val);
long long ans = query(0, n - 1, 1, l, r);
\end{lstlisting}


\subsection{Lazy Segment Tree}

Permite atualizações em intervalos sem visitar cada elemento.

\textbf{Construção:} $O(N)$.

\textbf{Consulta:} $O(\log N)$.

\textbf{Atualização em intervalo:} $O(\log N)$.

\begin{lstlisting}
int n;
vector<long long> tree;
vector<long long> lazy;

long long build(const vector<long long>& a, int l, int r, int p) {
    if (l == r)
        return tree[p] = a[l];

    int m = (l + r) >> 1;

    return tree[p] = build(a, l, m, p << 1) + build(a, m + 1, r, (p << 1) | 1);
}

void push(int l, int r, int p) {
    if (lazy[p] == 0)
        return;

    int m = (l + r) >> 1;

    // Aplica a atualizacao pendente nos dois filhos.
    tree[p << 1] += lazy[p] * (m - l + 1);
    tree[(p << 1) | 1] += lazy[p] * (r - m);

    lazy[p << 1] += lazy[p];
    lazy[(p << 1) | 1] += lazy[p];

    lazy[p] = 0;
}

long long query(int l, int r, int p, int ql, int qr) {
    if (r < ql || l > qr)
        return 0;

    if (ql <= l && r <= qr)
        return tree[p];

    push(l, r, p);

    int m = (l + r) >> 1;

    return query(l, m, p << 1, ql, qr) + query(m + 1, r, (p << 1) | 1, ql, qr);
}

long long update(int l, int r, int p, int ql, int qr, long long value) {
    if (r < ql || l > qr)
        return tree[p];

    if (ql <= l && r <= qr) {
        // Atualiza o intervalo inteiro sem descer na arvore.
        tree[p] += value * (r - l + 1);
        lazy[p] += value;

        return tree[p];
    }

    push(l, r, p);

    int m = (l + r) >> 1;

    return tree[p] = update(l, m, p << 1, ql, qr, value) + update(m + 1, r, (p << 1) | 1, ql, qr, value);
}
\end{lstlisting}

Inicialização:
\begin{lstlisting}
n = a.size();
tree.assign(4 * n, 0);
lazy.assign(4 * n, 0);
build(a, 0, n - 1, 1);
\end{lstlisting}

Uso:
\begin{lstlisting}
update(0, n - 1, 1, l, r, value);
long long answer = query(0, n - 1, 1, l, r);
\end{lstlisting}


\subsection{Merge Sort Tree}

Segment Tree onde cada nó mantém os elementos do seu
intervalo ordenados.

Permite consultas de ordem em intervalos usando
\texttt{lower\_bound} / \texttt{upper\_bound}.

\textbf{Construção:} $O(N\log N)$.

\textbf{Consulta:} $O(\log^2 N)$.

\begin{lstlisting}
int n;
vector<vector<int>> tree;

void build(const vector<int>& a, int p, int l, int r) {
    if (l == r) {
        tree[p] = {a[l]};
        return;
    }

    int m = (l + r) / 2;

    build(a, p * 2, l, m);
    build(a, p * 2 + 1, m + 1, r);

    // Cada no armazena os elementos do intervalo ordenados.
    merge(tree[p * 2].begin(), 
          tree[p * 2].end(), 
          tree[p * 2 + 1].begin(), 
          tree[p * 2 + 1].end(), 
          back_inserter(tree[p]));
}

int query(int ql, int qr, int x, int p, int l, int r) {
    // Intervalo atual nao participa da consulta.
    if (qr < l || r < ql)
        return 0;

    if (ql <= l && r <= qr)
        return lower_bound(tree[p].begin(), tree[p].end(), x) - tree[p].begin();

    int m = (l + r) / 2;

    return query(ql, qr, x, p * 2, l, m) + query(ql, qr, x, p * 2 + 1, m + 1, r);
}
\end{lstlisting}

Inicialização:
\begin{lstlisting}
n = a.size();
tree.assign(4 * n, {});
build(a, 1, 0, n - 1);
\end{lstlisting}

O exemplo acima conta elementos $<x$ no intervalo $[l,r]$.

Para contar elementos $\leq x$, use \texttt{upper\_bound}.


\subsection{DSU}

Disjoint Set Union / Union-Find.

Mantém componentes disjuntos e suporta união e consulta
de componentes.

\textbf{Complexidade com Compressão de Caminho:} $O(\alpha(N))$ por operação.

\textbf{Complexidade sem Compressão de Caminho:} $O(log(N))$ por operação.

\begin{lstlisting}
vector<int> parent;
vector<int> componentSize;

void init(int vertexCount) {
    parent.resize(vertexCount);
    componentSize.assign(vertexCount, 1);
    iota(parent.begin(), parent.end(), 0);
}

// Sem compressao de caminho.
int find(int vertex) {
    // Percorre a arvore ate encontrar a raiz.
    while (vertex != parent[vertex]) {
        vertex = parent[vertex];
    }

    return vertex;
}

// Com compresao de caminho.
int findCompressed(int vertex) {
    int root = vertex;

    // Primeiro encontra a raiz do conjunto.
    while (root != parent[root]) {
        root = parent[root];
    }

    // Faz todos os vertices do caminho apontarem para a raiz.
    while (vertex != root) {
        int nextVertex = parent[vertex];
        parent[vertex] = root;
        vertex = nextVertex;
    }

    return root;
}

bool unite(int firstVertex, int secondVertex) {
    int firstRoot = findCompressed(firstVertex);
    int secondRoot = findCompressed(secondVertex);

    if (firstRoot == secondRoot)
        return false;

    // Mantem a maior arvore como raiz.
    if (componentSize[firstRoot] < componentSize[secondRoot])
        swap(firstRoot, secondRoot);

    parent[secondRoot] = firstRoot;
    componentSize[firstRoot] += componentSize[secondRoot];

    return true;
}
\end{lstlisting}

Uso:
\begin{lstlisting}
init(vertexCount);
unite(firstVertex, secondVertex);
\end{lstlisting}